\documentclass{article}

\newcommand{\doctitle}{Assignment 6} % change document title
\usepackage{ellis_header}

\begin{document}

\section{Chapter 3.1}
\setcounter{problem}{1}
\begin{problem} % Problem 2
How many integers in $[100]$ are not divisible by 4, 6, or 7?
\end{problem}

\textbf{Answer:}
\begin{align*}
    N_{\ge}(\emptyset) = 100 \quad \quad & N_{\ge}(d_4, d_6) = \floor{\dfrac{100}{2 \cdot 6}} = 8 \\
    N_{\ge}(d_4) = \floor{\dfrac{100}{4}} = 25 \quad \quad & N_{\ge}(d_4, d_7) = \floor{\dfrac{100}{4 \cdot 7}} = 3 \\
    N_{\ge}(d_6) = \floor{\dfrac{100}{6}} = 16 \quad \quad & N_{\ge}(d_6, d_7) = \floor{\dfrac{100}{6 \cdot 7}} = 2 \\
    N_{\ge}(d_7) = \floor{\dfrac{100}{7}} = 14 \quad \quad & N_{\ge}(d_4, d_6, d_7) = \floor{\dfrac{100}{2 \cdot 6 \cdot 7}} = 1
\end{align*}

The answer is $100 - (25+16+14) + (8 + 3 + 2) - 1 = 57$

\setcounter{problem}{5}
\begin{problem} % Problem 6
    Answer the hat-check problem (i.e., the problem of derangement) for general $n$. This number is known as $D_n$. The question is ``How many ways can the hats of $n$ people be re-distributed so that each person receives exactly one hat, but no person receives there own hat''?
\end{problem}

\textbf{Answer:}

Before diving into the general $n$ solution. Lets look at an example party problem, with 3 people, called Kaelin, Alyse and Donna. The three cases to avoid are the following:
\begin{align*}
    h_k := & \: \text{``Kaelin gets his own hat back''} \\
    h_a := & \: \text{``Alyse gets her own hat back''} \\
    h_d := & \: \text{``Donna gets her own hat back''}
\end{align*}

So the set $P = \{h_k, h_a, h_d\}$ and the possible subsets of the problem are: $$J = \{\emptyset, \{h_k\}, \{h_a\}, \{h_d\}, \{h_k h_a\}, \{h_k h_d\}, \{h_a h_d\}, \{h_k h_a h_d\} \}$$ Instead of looping through each element in $J$ and adding them individually, a better approach would be to group cases where the same amount of people receive there own hat back. For instance.
\begin{itemize}
    \item $j = 0$, at least no one gets there hat back, $\emptyset$, 1 option 
    \item $j = 1$, at least one person gets there hat back, $\{h_k\}, \{h_a\}, \{h_d\}$, 3 options
    \item $j = 2$, at least two persons gets there hat back, $\{h_kh_a\}, \{h_kh_d\}, \{h_ah_d\}$, 3 options
    \item $j = 3$, at least three persons gets there hat back, $\{h_kh_ah_d\}$, 1 option
\end{itemize}
Each of the number of options can be represensented as ``With 3 people, how many ways can $j$ people get there hat back''? This is represented as $\binom{3}{j}$ or $\binom{n}{j}$ for a general amount of people, $n$. 

\medskip
Now we look at how to distribute the hats to people. If we didn't need to fix anyone getting there hat, the ways to distribute them would be $3!$. Kaelin would have 3 options, Alyse would have 2 options and Donna would have 1 option. If instead we fixed one person, say Kaelin, then you only have 2 options to choose from, $2!$. This is represented mathematically as $(3-j)!$ or $(n-j)!$ for the general case.

\medskip
Putting this together for the 3-person example we have:
\begin{equation*}
    D_3 = \sum_{j=0}^{3}\binom{3}{j}(-1)^{j}(3-j)!
\end{equation*}

Where the $(-1)^{j}$ is the book keeping for ensuring odd values are subtracted and evens added. This prevents accidental double counting of values that appear in multiple cases. For the general case, where its $n$ people, the equation is:
\begin{equation*}
    D_n = \sum_{j=0}^{n}\binom{n}{j}(-1)^{j}(n-j)!
\end{equation*}

Algebraically, this can be simplified starting with substiting the combination for its formula.

\begin{equation*}
    D_n = \sum_{j=0}^{n}\dfrac{n!}{j!(n-j)!}(-1)^{j}(n-j)! = n! \sum_{j=0}^{n}\dfrac{(-1)^{j}}{j!}
\end{equation*}

\begin{problem} % Problem 7
The number $D_n$ is the hat check problem for general $n$.
\begin{enumerate}[(a)]
    \item Calculate $\lim\limits_{n \rightarrow \infty} \dfrac{D_n}{n!}$
    \item Prove that for any $n, \: D_n$ equals the closest integer to $n!/e$.
\end{enumerate}
\end{problem}

\textbf{Answer:}

\begin{enumerate}[(a)]
    \item Looking at the limit of the hat check derangement over the total ways to distribute the hats.
    \begin{equation*}
        \dfrac{D_n}{n!} = \dfrac{n!}{n!} \sum_{j=0}^{n}\dfrac{(-1)^{j}}{j!} = \sum_{j=0}^{n}\dfrac{(-1)^{j}}{j!}
    \end{equation*}
    The inverse of Euler's number can be expressed with the following infinite series.
    \begin{equation*}
        \dfrac{1}{e} = \sum_{j=0}^{\infty}\dfrac{(-1)^{j}}{j!}
    \end{equation*}
    As $n \rightarrow \infty$ the summation approachs $1/e$. So the limit is the inverse of Euler's number.
    \item Explanation of the alternating series remainder term theorem from calculus \href{https://content.dodea.edu/VS/HS/AP_Calculus/websites/AP_Calculus_M9/section9_5c/docs/Lesson%2009.05c.06%20The%20Alternating%20Series%20Remainder%20Theorem.pdf}{here}.
    
    \medskip
    Before jumping right into using the alternating series remainder theorem, we need to first show the series converges. The criteria is:
    \begin{equation*}
        \lim_{j \rightarrow \infty}\dfrac{1}{j!} = 0 \quad \text{and} \quad \dfrac{1}{(j+1)!} < \dfrac{1}{j!}
    \end{equation*}
    So it passes. We can now apply the alternating series remainder theorem. This says the error of evaluated term will be less than or equal to the value of the first term left off. When $j=0 \text{ or } 1$ the series value is 1, 0. This doesn't make sense to begin evaluating the series at these values. The first value that makes sense is $j=2$, so two people.
    \begin{equation*}
        \sum_{j=0}^{2}\dfrac{(-1)^{j}}{j!} = 1 - 1 + 1/2 = 1/2
    \end{equation*}
    The next term is $j=3$ which gives $1/6$. So the low and high side of the $j=2$ case is $1/2 \pm 1/6 = (2/6, 4/6)$ or $(1/3, 2/3)$. The value this finally converges to is $1/e = 1.1/3$ which falls inside the boundary.

    \begin{proof}
    For any value of $n$ the number of derangements is:
    \begin{equation*}
        D_n = n! \sum_{j=0}^{n}\dfrac{(-1)^{j}}{j!}
    \end{equation*}
    The number of derangements over total permutations $D_n / n!$ is a converging series since it meets the alternating series theorem criteria.
    \begin{equation*}
        \lim_{j \rightarrow \infty}\dfrac{1}{j!} = 0 \quad \text{and} \quad \dfrac{1}{(j+1)!} < \dfrac{1}{j!}
    \end{equation*}
    We know that this series convergest to $1/e$ with $n$ approaching infinity.
    \begin{equation*}
        \dfrac{D_n}{n!} = \sum_{j=0}^{\infty}\dfrac{(-1)^{j}}{j!} = \dfrac{1}{e}
    \end{equation*}


    \end{proof}
\end{enumerate}

\setcounter{problem}{7}
\begin{problem} % Problem 8
How many ways can you distribute 20 identical objects to 10 distinct recipients so that each recipient receives at most five objects?
\end{problem}

\textbf{Answer:}

\begin{problem} % Problem 9
How many ways can you distribute $k$ identical objects to $n$ distinct recipients so that each recipient receives at most $r$ objects?
\end{problem}

\textbf{Answer:}

\setcounter{problem}{16}
\begin{problem} % Problem 17
A taxi drives from the intersection labeled $A$ to the intersection labeled $B$ in the grid of streets shown on page 94. The driver can only drive north (up) or east (right). Traffice reports indicate that there is heavy congestion at the intersections identified. How many routes from $A$ to $B$ can the driver take that \dots
\begin{enumerate}[(a)]
    \item Avoid all the congested intersections?
    \item Pass through at most one congested intersection?
\end{enumerate}
\end{problem}

\textbf{Answer:} Start by defining the different types of points. Ones that can go up and right. Ones that can only go up, ones that can only go right. Use the knowledge that each point has two choices (up or right) with a few exceptions.

\end{document}